\section{全市场排序}

31个一级行业中，7月6日至10日只有8个行业获得周度主力净流入。计算机净流入98.31亿元居首，但周涨2.74\%，行情已开始反应；银行净流入32.25亿元、周涨1.89\%，且处于52周区间约13\%位置，是最符合“资金先行、价格滞后”的行业。军工全周下跌2.12\%仍流入11.21亿元，本来具备潜伏特征，但7月10日单日资金集中涌入并推动板块上涨2.56\%，因此阶段升级为启动确认。

\begin{exhibitbox}[Exhibit 3：周度资金净流入行业]
\begin{tabularx}{\textwidth}{L{2.5cm}R{2.2cm}R{2.0cm}R{2.0cm}R{2.0cm}X}
\toprule
\textbf{行业} & \textbf{周资金/亿元} & \textbf{周涨跌} & \textbf{日资金/亿元} & \textbf{52周位置} & \textbf{判断} \\
\midrule
计算机 & 98.31 & +2.74\% & +5.21 & 15.5\% & 资金观察，已开始反应 \\
银行 & 32.25 & +1.89\% & +10.27 & 13.1\% & \textbf{静默吸纳首选} \\
传媒 & 27.33 & +2.01\% & +19.56 & 7.1\% & 早期启动确认 \\
国防军工 & 11.21 & -2.12\% & +80.49 & 30.5\% & 单日爆发，已非悄悄 \\
石油石化 & 8.81 & -1.04\% & -0.76 & 13.9\% & 周流入、当日转弱，观察 \\
煤炭 & 7.43 & +1.28\% & +0.07 & 43.4\% & 防御流入，弹性有限 \\
农林牧渔 & 2.15 & +0.33\% & +5.06 & 16.8\% & 低位试探，盈利分化 \\
房地产 & 1.35 & +1.89\% & +0.56 & 11.2\% & 低位但基本面门槛未过 \\
\bottomrule
\end{tabularx}
\sourcenote{证券时报数据宝；价格位置为AStock基于申万历史数据测算。}
\end{exhibitbox}

\section{单日反弹不能升级为布局}

\begin{exhibitbox}[Exhibit 4：7月10日流入、但全周仍流出的代表行业]
\begin{tabularx}{\textwidth}{L{2.7cm}R{2.2cm}R{2.2cm}R{2.0cm}X}
\toprule
\textbf{行业} & \textbf{当日资金/亿元} & \textbf{周资金/亿元} & \textbf{周涨跌} & \textbf{结论} \\
\midrule
医药生物 & +41.55 & -53.45 & -1.14\% & 政策催化后的单日回流，持续性未确认 \\
汽车 & +28.08 & -43.22 & -4.17\% & 低位修复，不是周度建仓 \\
公用事业 & +20.40 & -22.07 & -1.75\% & 防御承接，但全周仍失血 \\
机械设备 & +16.57 & -162.71 & -5.71\% & 行业整体流出，仅个别资金孤岛 \\
食品饮料 & +13.61 & -6.09 & +0.07\% & 估值低但尚无持续流入 \\
家用电器 & +7.07 & -31.08 & -1.92\% & 单日避险，未通过周度门槛 \\
\bottomrule
\end{tabularx}
\end{exhibitbox}

\section{完整31行业阶段表}

\small
\begin{longtable}{L{2.2cm}R{1.45cm}R{1.45cm}R{1.35cm}R{1.35cm}L{3.1cm}}
\toprule
\textbf{行业} & \textbf{周资金} & \textbf{周涨跌} & \textbf{52周位置} & \textbf{PE分位} & \textbf{阶段} \\
\midrule
\endfirsthead
\toprule
\textbf{行业} & \textbf{周资金} & \textbf{周涨跌} & \textbf{52周位置} & \textbf{PE分位} & \textbf{阶段} \\
\midrule
\endhead
银行 & +32.25 & +1.89\% & 13.1\% & 70.5\% & 静默吸纳（PB分位22.9\%） \\
国防军工 & +11.21 & -2.12\% & 30.5\% & 46.1\% & 启动确认 \\
煤炭 & +7.43 & +1.28\% & 43.4\% & 68.0\% & 防御型静默承接 \\
农林牧渔 & +2.15 & +0.33\% & 16.8\% & 40.5\% & 低位试探 \\
房地产 & +1.35 & +1.89\% & 11.2\% & 22.4\% & 低位但基本面未过门槛 \\
传媒 & +27.33 & +2.01\% & 7.1\% & 31.5\% & 早期启动 \\
计算机 & +98.31 & +2.74\% & 15.5\% & 41.7\% & 资金观察 \\
石油石化 & +8.81 & -1.04\% & 13.9\% & 33.6\% & 周流入、日转弱 \\
轻工制造 & -5.28 & -3.67\% & 2.3\% & 29.1\% & 单日反弹 \\
食品饮料 & -6.09 & +0.07\% & 5.5\% & 8.3\% & 单日反弹 \\
商贸零售 & -6.25 & -1.94\% & 3.9\% & 26.9\% & 单日反弹 \\
钢铁 & -6.54 & -4.12\% & 2.9\% & 52.6\% & 单日反弹 \\
纺织服饰 & -6.71 & -4.56\% & 5.8\% & 30.9\% & 单日反弹 \\
环保 & -17.84 & -4.88\% & 25.0\% & 25.4\% & 单日反弹 \\
交通运输 & -21.69 & -1.46\% & 7.6\% & 34.1\% & 单日反弹 \\
公用事业 & -22.07 & -1.75\% & 16.5\% & 42.2\% & 单日反弹 \\
家用电器 & -31.08 & -1.92\% & 12.4\% & 30.2\% & 单日反弹 \\
汽车 & -43.22 & -4.17\% & 10.0\% & 59.7\% & 单日反弹 \\
医药生物 & -53.45 & -1.14\% & 27.0\% & 27.2\% & 单日反弹 \\
机械设备 & -162.71 & -5.71\% & 81.0\% & 79.2\% & 单日反弹/内部孤岛 \\
有色金属 & -253.80 & -7.66\% & 45.8\% & 23.5\% & 单日反弹 \\
美容护理 & -1.23 & -2.93\% & 7.2\% & 13.4\% & 低位无资金 \\
社会服务 & -11.19 & -4.68\% & 2.0\% & 20.1\% & 低位无资金 \\
建筑装饰 & -46.58 & -3.10\% & 13.1\% & 65.0\% & 低位无资金 \\
非银金融 & -112.48 & -2.54\% & 16.6\% & 1.2\% & 低位无资金 \\
综合 & -2.62 & -2.07\% & 60.9\% & 96.9\% & 无信号 \\
通信 & -83.21 & -1.66\% & 75.9\% & 46.9\% & 高位兑现 \\
建筑材料 & -91.03 & -12.43\% & 63.6\% & 99.4\% & 下跌但估值不低 \\
基础化工 & -211.10 & -8.22\% & 49.3\% & 58.7\% & 周度失血 \\
电力设备 & -395.10 & -8.87\% & 48.1\% & 41.6\% & 周度大幅失血 \\
电子 & -489.47 & -1.66\% & 82.2\% & 99.4\% & 高位拥挤/内部强分化 \\
\bottomrule
\end{longtable}
\normalsize

银行的PE分位并不低，原因是盈利增长慢、市场给予稳定性溢价；真正体现安全边际的是PB分位与高股息。因此银行必须用PB--ROE和资产质量，而不是机械使用PE判断“便宜”。
